\chapter{Frozen Evaluation Artifact}
\label{app:frozen}
% THE AUDIT TRAIL. Everything frozen before evaluation day, with git commit
% hashes and dates:
%   - the Recovery-Outcome Rubric levels and thresholds
%   - the scoring script version
%   - the workflow baseline tag (Section~\ref{sec:wf:freeze})
%   - the parity conformance test result
% This appendix is what makes the preregistration claim checkable rather than
% merely stated.

\paragraph{Baseline freeze.} The workflow baseline was frozen in commit
\texttt{67a14f8}, made at 11:41 on 2026-09-10 and tagged
\code{eval-day-2026-09-10}. The five novel scenarios were saved between 15:27
and 16:04 the same day and then committed.

\paragraph{Pre-registration tags.} The tags \code{prereg-v1} and
\code{prereg-amendment-1} were cut together on 2026-09-12, on commit
\texttt{7a4d7cc}. This is the commit the anticipated campaign ran on. The pilot
ran before the tag, as the pre-registration discloses.

\paragraph{The one post-reveal change.} Commit \texttt{e9e20cd}, merged after
the reveal, added the switch that turns off re-planning for the no-replan
control (Amendment~6). It also merged two fenced sections into
\code{thresholds.py}, one for the negative control and one for the decision
rubric. Neither section holds anything the ladder reads. The threshold hash
recorded with the scored runs, \texttt{d08c1e35}, is identical across all 310
of them.

\chapter{Disruption Brief}
\label{app:brief}
% The characteristics-only brief given to the external supplier, reproduced in
% full. Shows what was and was not communicated, which is the firewall evidence.

\chapter{Adjudication Log}
\label{app:adjudication}
% Written record of every residual edge call the deterministic scorer could not
% settle (mainly level-0 vs level-1), with the reasoning for each.

\chapter{Database Schema (DDL)}
\label{app:ddl}
% TODO: \lstinputlisting[language=SQL]{...}

\chapter{Agent Prompts}
\label{app:prompts}
% TODO: CORE system prompt + the three event playbooks, per agent.

\chapter{Scenario Definitions}
\label{app:scenarios}
% TODO: \lstinputlisting of the scenario JSON files.

\chapter{Decision Rubric Counts}
\label{app:rubric}

Table~\ref{tab:res:rubric} gives the counts behind
Figure~\ref{fig:res:rubric}, grouped by disruption.
Table~\ref{tab:eval:heldout} states what each criterion checks. Each cell
counts \texttt{yes} judgements out of 10.

\begin{table}[htbp]
\centering
\caption{Held-out decision rubric, \texttt{yes} judgements out of 10.}
\label{tab:res:rubric}
\footnotesize
\begin{tabular}{lcccc}
\toprule
Criterion & baseline & qwen3.5 & glm-5.2 & deepseek \\
\midrule
\multicolumn{5}{l}{\emph{Disruption:} \code{bus1_suspension_capacity_reduction}} \\
\quad \code{reserve_committed} & 10 & \textbf{10} & \textbf{10} & 9 \\
\quad \code{whole_population_planned} & 10 & \textbf{10} & \textbf{10} & 6 \\
\quad \code{bus_load_within_cap} & 0 & 0 & 3 & 2 \\
\midrule
\multicolumn{5}{l}{\emph{Disruption:} \code{conflicting_evacuation_orders}} \\
\quad \code{authority_informed} & 10 & \textbf{10} & 9 & \textbf{10} \\
\quad \code{full_evacuation_kept} & 10 & 9 & \textbf{10} & 5 \\
\quad \code{facility_answered_in_time} & \textbf{10} & 0 & 0 & 0 \\
\midrule
\multicolumn{5}{l}{\emph{Disruption:} \code{medical_emergency_in_flight}} \\
\quad \code{authority_informed} & 10 & \textbf{10} & \textbf{10} & \textbf{10} \\
\quad \code{in_flight_leg_lands_where_ordered} & 10 & 8 & 6 & 6 \\
\quad \code{medic_sent_to_the_bus_destination} & 0 & 1 & 1 & 0 \\
\midrule
\multicolumn{5}{l}{\emph{Disruption:} \code{unverified_persons_report}} \\
\quad \code{verification_requested} & 10 & 9 & \textbf{10} & \textbf{10} \\
\quad \code{stood_down_after_retraction} & 10 & \textbf{10} & \textbf{10} & 8 \\
\quad \code{reserve_committed_in_time} & 10 & 9 & \textbf{10} & 9 \\
\quad \code{facility_runs_untouched} & \textbf{10} & 7 & 8 & 5 \\
\midrule
\multicolumn{5}{l}{\emph{Disruption:} \code{sc1_partial_capacity_loss}} \\
\quad \code{no_order_raises_shelter_above_beds} & \textbf{10} & 5 & 9 & 8 \\
\quad \code{shelter_set_down_within_beds} & 0 & 2 & 2 & 0 \\
\quad \code{everyone_bedded_at_the_reduced_capacity} & 0 & 0 & 0 & 0 \\
\bottomrule
\end{tabular}
\end{table}

\chapter{Reproducibility}
\label{app:repro}
% How to run the systems and the harness. Model digests, environment variables
% (SIM_TIME_SCALE, OLLAMA_MODEL, ORS_API_KEY), migrations, seeds.

\paragraph{The headless world client.} On the agentic side, a bus leg closes
only when the world acknowledges that the task is complete. In the first pilot
that acknowledgement came only from the browser's vehicle animation, so without
a browser no bus ever arrived and every agentic cell scored level 0. We ruled
that pilot invalid and re-ran it on the same plan. Every agentic cell runs a
headless client that plays the animations from the backend's own events. A
cell refuses to start while a browser tab is connected, since both would
acknowledge every leg.

\paragraph{Replicate blocks.} The replicate index shifts the policy seed, so
blocks are kept apart by index and by output file. Replicates 100--105 are the
pilot, 1000--1009 the confirmatory runs, and 2000 onwards the model
screenings. The scoring and summary commands read only the confirmatory file
by default, so the code itself keeps pilot and screening runs out of the
confirmatory analysis.

\paragraph{Two-pass scoring.} A cell's clean-run reference $K^\ast$ must exist
before any run in it may rise to level 4. Scoring therefore runs in two
passes, first the references and then the levels.

\paragraph{Wilcoxon variants.} The signed-rank test is exact for up to 20
non-zero differences without ties. Otherwise it uses a normal approximation
with tie correction.
